% generated from papers/XXXIV-bottleneck/paper_XXXIV_bottleneck.tex -- edit the paper, then rerun assemble.py
\chapter[The abelian scaling bottleneck and the angular system]{The abelian scaling bottleneck, and the angular Hurwitz system:\\ where noncommutative arithmetic thermodynamics can and cannot live}\label{ch:XXXIV}
\renewcommand{\Prob}{\mathrm{Prob}}
\chapterorigin{This chapter is Paper XXXIV of the series (\texttt{papers/XXXIV-bottleneck/}).}
\begin{summary}
The series has twice failed to produce a genuinely noncommutative arithmetic
thermodynamics: the free semigroup (Chapter~\ref{ch:V}) loses the high-temperature phase to
orthogonal range projections, and the Hurwitz quaternions (Papers X, XII) have nested
ranges but a scalar dynamics driven by the reduced norm, which factors through the
abelianization. We prove that the second failure is universal --- the \emph{abelian
scaling bottleneck}: for a right LCM semigroup with scalar one-parameter dynamics
$\sigma_t(v_a)=N(a)^{it}v_a$, the KMS classification data splits into a scaling measure
governed by the $P^{\mathrm{ab}}$-graded weight $N$ and a $\beta$-independent tracial datum on
the isotropy; every $\beta$-driven transition is a transition of the abelianized weight
data, and no nonabelian residual symmetry group can be produced on the semigroup side
--- consistently with Chapter~\ref{ch:VII}'s finding that its essential-value group is abelian, and
verified concretely on the braid monoids $B_n^+$, right LCM and maximally noncommutative,
whose entire transition is that of the length series while the braiding survives only as
$\beta$-independent traces. The complementary half is Chapter~\ref{ch:XXVII}: noncommutativity lives
\emph{transverse} to the scaling, as a cocycle into a compact group, realizing
$\Prob(H\backslash G)$. We then name the Goldilocks candidate the two halves point to:
the \emph{angular Hurwitz system}, whose radial part $\Nrd$ carries the thermodynamics of
Chapter~\ref{ch:X} ($\beta_c=2$) and whose discarded angular part $\rho(a)=a/\Nrd(a)^{1/2}\in SU(2)$
is an ordered cocycle in the sense of Chapter~\ref{ch:XXVII}. The $\pm a$ symmetry of the
representation sets makes the matrix Kakutani series exactly computable,
$\Xi_\beta(\rho)=96\sum_p(p{+}1)p^{-\beta}$ up to the projective normalization, with
abscissa \emph{exactly} the radial critical exponent $2$: the angle turns critical
precisely where the radius does. The programme --- transport of Chapter~\ref{ch:XXVII}'s matrix
martingale to the Hurwitz tail relation in the Gibbs phase, expected verdict
$\mathrm{KMS}_\beta$-simplex $\cong\Prob(SU(2))$ for $\beta>2$, with the
Lubotzky--Phillips--Sarnak spectral gap closing the other side --- is stated with its
single hard step named.
\end{summary}


\section{The bottleneck theorem}\label{XXXIV:sec:bottleneck}

Let $P$ be a countable right LCM monoid with trivial units for simplicity, $N:P\to
\R_{>0}$ a multiplicative weight with finite level sets, $\mathcal T(P)$ its Nica--Toeplitz
algebra with generating isometries $v_a$, nested range projections, and the scalar
dynamics $\sigma_t(v_a)=N(a)^{it}v_a$.

\begin{lemma}[The trivial half, stated because it is the point]\label{XXXIV:lem:ab}
$\R_{>0}$ is abelian, so $N$ factors through the abelianization $P^{\mathrm{ab}}$: the Gibbs
weight of every monomial $v_av_b^*$ depends only on the classes of $a,b$ in $P^{\mathrm{ab}}$.
\end{lemma}

\begin{theorem}[Abelian scaling bottleneck]\label{XXXIV:thm:bottleneck}
Let $\varphi$ be a $\sigma$-KMS$_\beta$ state on $\mathcal T(P)$ (or any gauge-invariant
quotient). Then:
\begin{enumerate}
\item On the fixed-point algebra of the gauge action, spanned by $v_av_b^*$ with
$N(a)=N(b)$, $\varphi$ is a trace; in particular on every subalgebra generated by
$N$-preserving symmetries (the ``kernel directions'') the state is tracial, with no
$\beta$-dependence in the defining condition.
\item $\varphi$ is determined by the pair (its restriction to the diagonal, a scaling
measure with cocycle $N^{-\beta}$; its tracial field on the isotropy), in the
Neshveyev parametrization \cite{Nesh13}; the scaling measure sees only the
$P^{\mathrm{ab}}$-graded counting data of $N$.
\item Consequently the $\beta$-driven structure of the KMS simplex --- existence ranges,
uniqueness, symmetry breaking --- is that of an abelianized weight system, and the
residual symmetry acting on extreme states through the scaling-measure component is a
quotient of a group of characters: no scalar dynamics on a right LCM semigroup produces
a nonabelian residual symmetry group on that component.
\end{enumerate}
\end{theorem}

\begin{proof}[Proof sketch]
(1) is the two-line computation used throughout the series (Chapter~\ref{ch:XXVII}, its Theorem
2.1): for $x$ in the fixed algebra and analytic $y$ there, KMS gives
$\varphi(xy)=\varphi(y\sigma_{i\beta}(x))=\varphi(yx)$ since $\sigma$ is trivial on the
fixed algebra. (2) is the groupoid form of the KMS condition: the diagonal restriction is
quasi-invariant with Radon--Nikodym cocycle $N^{-\beta}$, whose value on any groupoid
element depends only on the $P^{\mathrm{ab}}$-class by Lemma~\ref{XXXIV:lem:ab}, and the fibre datum
over the isotropy is a field of tracial states, weightless because the cocycle is trivial
on isotropy of $N$-preserving type; extremality decomposes accordingly \cite{Nesh13}.
(3) follows: two extreme states with the same measure component differ by a
$\beta$-independent tracial datum, and the action of any symmetry on the measure
component is through the characters detected by the $N^{-\beta}$-weighted series of
[Main], an abelian dual object.
\end{proof}

\begin{example}[Braid monoids: maximal noncommutativity, abelian transition]\label{XXXIV:ex:braid}
The braid monoid $B_n^+$ is right LCM (Garside theory \cite{Dehornoy}) and strictly
noncommutative for $n\ge3$; its abelianization is $\Z$, generated by word length
$\ell$, and any scalar dynamics is $\sigma_t(v_a)=q^{it\ell(a)}$ for some $q>1$. By
Theorem~\ref{XXXIV:thm:bottleneck} the entire KMS transition structure is that of the growth
series $\sum_aq^{-\beta\ell(a)}$ --- a radius-of-convergence phenomenon on $\Z$, with at
most circle-character symmetry --- while the braiding survives only in the
$\beta$-independent tracial fibre. The Goldilocks checklist (noncommutative, nested,
coupled) fails at the third item, and the theorem says it must.
\end{example}

\begin{remark}[Consistency across the series]\label{XXXIV:rem:consistency}
Theorem~\ref{XXXIV:thm:bottleneck} retro-explains three findings: Chapter~\ref{ch:VII}'s essential-value
group $E_\beta$ came out abelian because it had to; Chapter~\ref{ch:XII}'s reduction of the search
to the norm-one subgroup located the noncommutativity exactly in the weightless kernel
directions, where item (1) makes it tracial; and Chapter~\ref{ch:XXVIII}'s no-go then showed those
tracial directions are, for arithmetic orders, hyperfinite. The bottleneck is the common
cause.
\end{remark}

\section{The complementary half, and where the grail lives}\label{XXXIV:sec:complement}

The bottleneck forbids noncommutativity in the \emph{scaling} direction. Chapter~\ref{ch:XXVII}
proved it thrives \emph{transverse} to it: keeping the scalar dynamics and the principal
tail relation, an ordered Frobenius cocycle $c$ into a nonabelian compact group $G$
yields KMS simplices $\Prob(H\backslash G)$, $H$ the Mackey range, with genuinely
nonabelian breaking, lift-dependence, and the matrix Kakutani criterion (its Theorem
4.3). The two statements are the two halves of one verdict: \emph{scaling abelian,
symmetry arbitrary}. The grail, if it exists, is a single arithmetic semigroup whose
canonical structure produces such a cocycle intrinsically.

\section{The angular Hurwitz system}\label{XXXIV:sec:angular}

The Hurwitz monoid $P$ of Chapter~\ref{ch:X} discards, in its dynamics, exactly half of each
generator: $a=\Nrd(a)^{1/2}\cdot\rho(a)$ with
\[
\rho(a):=a/\Nrd(a)^{1/2}\ \in\ SU(2)\quad(\text{unit quaternions}),
\]
the \emph{angular part}, multiplicative up to the ordering issues that Chapter~\ref{ch:XXVII}'s
construction is built to handle: fixing the norm-ordering of the primes and, within a
level, the self-conjugate generator sets of Chapter~\ref{ch:XXVIII}, the assignment
$a\mapsto\rho(a)$ defines an ordered cocycle $c_\rho$ on the Hurwitz tail relation with
values in the projective unit group $SU(2)/\{\pm1\}\cong SO(3)$, and lifts to $SU(2)$
level by level.

\begin{proposition}[The angular Kakutani series, exactly]\label{XXXIV:prop:angular}
For the level-$p$ generator set $R_p$, of cardinality $24(p{+}1)$ and closed under
$a\mapsto-a$ and $a\mapsto\bar a$, the Sato--Tate average is exact:
$\sum_{a\in R_p}\mathrm{tr}\,\rho(a)=\frac{2}{\sqrt p}\sum_{a\in R_p}x_0(a)=0$, by the
symmetry $a\mapsto-a$. Hence
\[
\Xi_\beta(\rho)\ =\ \sum_p p^{-\beta}\sum_{a\in R_p}\bigl\|I-\rho(a)\bigr\|_{\HS}^2
\ =\ \sum_p p^{-\beta}\cdot 4\,|R_p|
\ =\ 96\sum_p(p{+}1)\,p^{-\beta},
\]
with abscissa of convergence exactly $2$ --- the radial critical exponent
$\beta_c(\zeta_P)=2$ of Chapter~\ref{ch:X}. The angle turns critical precisely where the radius
does. (Projectively, replacing $\rho$ by its $SO(3)$-image replaces the constant $4$ by
the corresponding adjoint-trace constant; the abscissa is unchanged.)
\end{proposition}

\begin{proof}
$\|I-u\|^2_{\HS}=4-2\,\mathrm{Re}\,\mathrm{tr}\,u$ on $SU(2)$; sum over $R_p$ and use
$-a\in R_p$, which cancels the trace term identically; $|R_p|=24(p{+}1)$ is the
Hurwitz--Jacobi count. Convergence of $\sum(p{+}1)p^{-\beta}$ is $\beta>2$.
\end{proof}

\begin{remark}[The programme, with its hard step named]\label{XXXIV:rem:programme}
In the Gibbs phase $\beta>2$ of Chapter~\ref{ch:X}, $\Xi_\beta(\rho)<\infty$; if Chapter~\ref{ch:XXVII}'s
matrix-martingale coboundary theorem (its Theorem 4.3) transports from the independent
geometric relation to the Hurwitz tail relation --- the single hard step: the
coordinates of the Hurwitz boundary are neither independent nor geometric, and the
martingale filtration must follow the tree levels of Chapter~\ref{ch:XXVIII} --- then $c_\rho$ is a
coboundary there and the KMS$_\beta$ simplex of the angular system is
$\Prob(SU(2))$: a simplex over a genuinely nonabelian compact group, carried by an
honest arithmetic semigroup, in its entire low-temperature phase. On the other side, the
Lubotzky--Phillips--Sarnak spectral gap \cite{LPS} makes every nontrivial representation
of the level-$p$ generators uniformly contractive, which is exactly the mixing hypothesis
of Chapter~\ref{ch:XXVII}'s Theorem 4.5; where a scaling measure exists at the critical edge, the
angular symmetry should therefore be fully restored. Verdict: the bottleneck theorem
closes the scaling route universally; the angular route is open, constructive, and
reduced to one named transport problem.
\end{remark}

